\chapter*{\hfill{\centering  ЗАКЛЮЧЕНИЕ}\hfill}
\addcontentsline{toc}{chapter}{ЗАКЛЮЧЕНИЕ}

В результате исследования было установлено, что:
\begin{itemize}[label=---]
	\item линейный алгоритм поиска расстояния Левенштейна работает быстрее линейного алгоритма поиска расстояния Дамерау-Левенштейна в $\sim$1.235 раз;
	\item рекурсивная реализация алгоритма Левенштейнва с кешированием работает значительно быстрее (в $\sim$299.15 раз), чем без кеширования;
	\item рекурсивная реализация алгоритма линейного поиска с кешированием медленнее линейной реализаци в $\sim$2.36 раз.
\end{itemize}
Также было установлено, что рекурсивные алгоритмы по памяти более выгодны, чем линейные.

Цель данной лабораторной работы: описание и исследование алгоритмов поиска расстояний Левенштейна и Дамерау---Левенштейна, была достигнута.

При достижении поставленной цели были выполнены следующие задачи:
\begin{enumerate}
	\item Рассмотрены следующие алгоритмы посика:
	\begin{enumerate}
		\item линейный алгоритм поиска расстояния Левенштейна;
		\item рекурсивный алгоритм поиска расстояния Левенштейна;
		\item рекурсивный алгоритм поиска расстояния Левенштейна c кешированием;
		\item линейный алгоритм поиска расстояния Дамерау---Левенштейна.
	\end{enumerate}
	\item Написана программа, реализующая вышеперечисленные алгоритмы поиска.
	\item Проведено исследование затрачиваемого процессорного времени и памяти при различных реализациях алгоритмов.
	\item Проведён сравнительный анализ алгоритмов.
\end{enumerate}
